

\section{The Standard IS-LM Model}

\begin{frame}{IS-LM Model: IS Curve}
    \showmytikz[1,2,2.5]{../shared/tikz/tikz-islm}
    \mycomment{
   $IS$ is the curve of all interest rate - output - combinations that result in a balance of the real markets, that is, of investments~$I$ and savings~$S$.

    At $P_1$, a relatively high interest rate makes credit expensive. At this rate, much will be saved, hence less will be spent on consumption. There will be relatively low output.

    At $P_2$, lower interest rate make consumption more interesting than saving. Thus, output $Y$ will be higher.

    % Empirically, the $IS$ curve is relatively steep, while the $LM$ curve is relatively flat. 
    }
\end{frame}

\begin{frame}{IS-LM Model: LM Curve}
    \showmytikz[3,4,4.5]{../shared/tikz/tikz-islm}
    \mycomment{
        $LM$ is the curve of all combinations of interest rate and output that result in a balance on the money market, namely of liquidity preference $L$ and money supply $M$.

        At $P_3$, a relatively low output will require little money, and low money demand will result in low interest rates. 

        At $P_4$, a higher output will require more money, and as more profits are realised, interest rates will rise.

        The exact shape of the $LM$ curve is less clear than that of the $IS$ curve. In reality, it might be very flat. Some say it is flat on the left (liquidity trap) and steep on the right.        
        
        % Future expectations and policy play a larger role on the financial markets than in the real markets. 
    }
\end{frame}

\begin{frame}{IS-LM Model: Equilibrium}
    \showmytikz[1,3,5]{../shared/tikz/tikz-islm}
    \mycomment{
        There is an equilibrium $A$ where $IS$ and $LM$ meet, that is the synchronous balance of real and financial markets. At this combination of interest rate $i$ and output $Y$ investments will equal savings (the point is on $IS$) and money supply will equal money demand (the point is on $LM)$.
    }
\end{frame}


\begin{frame}{IS-LM Model: Disequilibria}
    \showmytikz[1,2,3,4,5]{../shared/tikz/tikz-islm}
    \mycomment{
        $P_1$: high interest rates result in little spending and low output. At this interest rate there is an excess of money supply (the point is left of $LM$). Money suppliers will lower the interest rate to find clients, which will result in a higher output, so movement towards $A$.

        $P_2$: at a low interest rate and high output, there is excess demand for money for transactions. This will drive up interest rates. 

        $P_3$: low interest rate and low output means there is much cheap money for more consumption available, creating excess demand for goods, driving up output. 

        $P_4$: high interest rates clear the money market, but the output is too high at this point for a balanced real market; subjects would rather save than consume this much. Thus, production and interest rates will have to go down.
        
    }
\end{frame}



\begin{frame}{IS-LM Model: Fiscal Policy}
    \showmytikz[1,3,5,6]{../shared/tikz/tikz-islm}
    \mycomment{
        Government expenditure can shift the $IS$ curve to the right, resulting in higher interest rate $i_2>i_1$ and larger output $Y_2>Y_1$. 
        
        The effect is only partial; the rise in output is smaller than the additional government expenditure. The risen interest rate has resulted a partial crowding-out. 

        Tariffs have a similar effect, as they raise net exports ($NX=EX-IM$). 

        Any fiscal policy has also distributive effects not shown here.
    }
\end{frame}

\begin{frame}{IS-LM Model: Monetary Policy}
    \showmytikz[1,3,5,7]{../shared/tikz/tikz-islm}
    \mycomment{
        Monetary policy (by the central bank) can shift the $LM$ curve to the right. This can happen either through lowering of the interest rate, or by open market policy. 

        The new equilibrium shows a lower interest rate and higher output $Y_3>Y_1$. The effect is also only partial: the rise in output is smaller than the monetary stimulus, since the lower interest rate has reduced savings and investment.

        There is an obvious contradiction (or: competition?) between fiscal and monetary policy, which is the reason why central banks \textit{should be} independent.
    }
\end{frame}

\begin{frame}{IS-LM Model: Fiscal vs. Monetary Policy}
    \showmytikz[1,3,5,6,7]{../shared/tikz/tikz-islm}
    \mycomment{
        Fiscal and monetary policy can complement or offset each other. Expansion by policy will meet its natural boundary at unsustainable government debt levels or at runaway inflation discrediting the currency. 
        
        It is thus widely believed that any such policy should be used temporarily only, like to overcome an acute crisis, external shocks, slumps, or maybe to catch up on development. 
        
        It is a bit like burning furniture to heat the room, or selling the silverware: you can do it in a crisis, but it is not a long-term solution.
    }
\end{frame}


   
\begin{frame}{IS-LM Model, Curved: Fiscal Policy}
    \showmytikz[1,3,5,6]{../shared/tikz/tikz-islm-smooth}
    \mycomment{
        The actual shape of IS and LM are debated. They `seem' to be steep at the extremes, but flat for low levels of $i$. This can, of course, have dramatic consequences whether fiscal or monetary policy works at all. 

        A steep $IS$-curve indicates that fiscal policy has only little impact on interest rates. Then there is relatively little crowding-out.
    }
\end{frame}


    \begin{frame}{IS-LM Model, Curved: Monetary Policy}
    \showmytikz[1,3,5,7]{../shared/tikz/tikz-islm-smooth}
    \mycomment{
        The shape of the $LM$-curve is much disputed. The consensus sees it flat for low $Y,i$ and steep for higher $Y,i$. 

        In the lower region, the flat $LM$-curve implies effectivelessness of monetary policy; larger money supply does not lead to more investment, since economic subjects have a strong preference for liquidity over investing. This situation is called `liquidity trap'.
    }
\end{frame}


    
